\section{Geometric realization of two marked collision profiles}
\label{mapsD:repair:modular-collisions}

We use positive base meridians and column vectors for integral elliptic
homology.  A positive nodal monodromy is
\[
 T=\begin{pmatrix}1&1\\0&1\end{pmatrix},\qquad
 S=\begin{pmatrix}0&-1\\1&0\end{pmatrix},\qquad
 R=ST=\begin{pmatrix}0&-1\\1&1\end{pmatrix}.
\]
Thus $s=[S]$ and $r=[R]$ generate
$G=\operatorname{PSL}_2(\mathbb Z)=\langle s,r\mid s^2=r^3=1\rangle$,
and $[T]=sr$.  Products of matrices are written in the order used for
the archived cocycles: $A_1\cdots A_n=I$.  The section cocycle
$(q_i)$ satisfies
\begin{equation}\label{mapsD:repair:cocycle-relation}
 q_i\in\operatorname{sat}\operatorname{im}(A_i-I),\qquad
 \sum_{i=1}^n P_iq_i=0,\qquad P_i=A_1\cdots A_{i-1}.
\end{equation}
Changing one global logarithm by $v\in\mathbb Z^2$ changes $q_i$ by
$(A_i-I)v$.  Saturation here means intersection of the rational image
with $\mathbb Z^2$.

The fibre configurations and the rank-one Mordell--Weil lattices in
the next theorem are contained in Miranda's existence list and
Fukae's Table~3; they are not asserted to be new configurations
\cite[\S4, pp.~206--210]{manMiranda90}
\cite[Table~3, rows~50 and~56]{manFukae02}.
The additional verification is the realization of the particular
based monodromy and section cocycles used in the calculations.

\begin{theorem}[Realization of the archived marked profiles]
\label{mapsD:repair:collision-realization}
\StatementRecord{mapsD:repair:collision-realization}{proved}{Historical events: GPT-6 Astra (OpenAI), ChatGPT Pro mode, effort Pro, S6\_09\_PAPERS: carried as background without fault found. GPT-6 Astra (OpenAI), ChatGPT Pro mode, effort Pro, S6\_09\_PAPERS, 17 November 2026: targeted proof check. Also checked in this writing. Pre-continuation writing events: GPT-6 Astra (OpenAI), ChatGPT Pro mode, effort Pro, S6\_10\_PAPERS, 18 November 2026: targeted proof check. Continuation metadata events: GPT-6 Astra (OpenAI), ChatGPT Pro mode, effort Pro, S6\_10\_PAPERS continuation, 19 November 2026: used and rechecked in this run; status-vocabulary reconciliation; metadata reconciliation only; no new mathematical check. \emph{Pre-continuation recorded proof scope.} GPT-6 Astra (OpenAI, ChatGPT Pro mode; recorded reviewer /root/cdp\_audit), effort Pro, 17 November 2026; GPT-6 Astra (OpenAI, ChatGPT Pro mode), effort Pro, 2026-11-17 UTC; also checked in this writing by GPT-6 Astra (OpenAI), ChatGPT Pro mode, effort Pro, 18 November 2026}{the hypotheses, notation and definitions stated here; \cite[pp.~35--38]{mapsMilne2017}; \cref{mapsD:repair:50-weierstrass}; \cite[\S8.2--8.4, pp.~38--39]{manSS10}; \cref{mapsD:repair:cocycle-relation}; \cite[Lemma~4.1, Theorem~4.2 and Corollaries~4.2--4.3,
pp.~21--22]{manFukae02}; \cite[Lemmas~3.1--3.2, (11), pp.~11--13]{manFukae02}}

There are rational elliptic surfaces with zero section over
$\mathbb P^1$ having, in the indicated order, the following singular
fibres and positive based monodromy matrices:
\[
\begin{array}{c|c|c}
\text{profile}&\text{fibre types}&(A_1,\ldots,A_n)\\ \hline
50/\mathrm{II}&(\mathrm{II},I_3,I_1^*)&
\left(
\begin{pmatrix}3&7\\-1&-2\end{pmatrix},
\begin{pmatrix}7&12\\-3&-5\end{pmatrix},
\begin{pmatrix}-2&-1\\1&0\end{pmatrix}\right)\\[5pt]
56/\mathrm{II}&(I_2,\mathrm{II},I_5,I_3)&
\left(
\begin{pmatrix}5&8\\-2&-3\end{pmatrix},
\begin{pmatrix}27&37\\-19&-26\end{pmatrix},
\begin{pmatrix}61&80\\-45&-59\end{pmatrix},
\begin{pmatrix}4&3\\-3&-2\end{pmatrix}\right).
\end{array}
\]
Their Mordell--Weil groups are infinite cyclic.  A generator can be
chosen with the respective integral logarithm cocycle
\[
 q^{50}=\bigl((-5,2),(-8,4),(0,1)\bigr),\qquad
 q^{56}=\bigl((-2,1),(0,0),(4,-3),(1,-1)\bigr).
\]
Its height is $1/12$ in profile $50/\mathrm{II}$ and $1/30$ in
profile $56/\mathrm{II}$.  These are the exact marked data in the
archived collision-profile file, not merely data with the same fibre
types or the same height.
\end{theorem}


\noindent\textsf{Uses:} the finite-index modular-curve construction
\cite[Propositions~2.16, 2.21; Theorem~2.22; Proposition~4.3 and
Remark~4.4]{mapsMilne2017}; the minimal Weierstrass criterion
\cite[\S4.5, Table~1; \S8.2--8.4]{manSS10}; and the junction
identification \cite[\S3.1, (4), (8), (9); Lemmas~3.1--3.2, 4.1;
Corollaries~4.2--4.3]{manFukae02}.

\begin{proof}
\begin{enumerate}[label=\arabic*.,leftmargin=*]
\item \par\medskip\noindent\textbf{1. Two specified modular covers.}\quad 
We first construct the covers, retaining their markings.  A row of
either of the following tables gives the right action of $s,r,r^{-1}$
on a coset labelled by its first entry:
\[
\begin{array}{c|rrr}
50/\mathrm{II}&s&r&r^{-1}\\\hline
0&1&2&3\\1&0&1&1\\2&3&3&0\\3&2&0&2
\end{array}
\qquad
\begin{array}{c|rrr}
56/\mathrm{II}&s&r&r^{-1}\\\hline
0&1&2&3\\1&0&1&1\\2&4&3&0\\3&5&0&2\\
4&2&6&7\\5&3&8&9\\6&9&7&4\\7&8&4&6\\
8&7&9&5\\9&6&5&8
\end{array}
\]
Each $s$-permutation is an involution, each $r$-permutation has cube
one, and the action is transitive.  These statements can be checked
directly in the displayed tables.  Consequently the stabilizer of
the coset $0$ is a subgroup of index $4$, respectively $10$, in $G$.
Its inverse image in $\operatorname{SL}_2(\mathbb Z)$ contains $-I$.
We use the compact modular curve of this subgroup.  The construction
of its complex structure uses the coordinates $z^e$ at elliptic
points and $\exp(2\pi i\tau/w)$ at a cusp of width $w$;
it applies to every finite-index subgroup and does not require a
congruence assumption \cite[pp.~35--38]{mapsMilne2017}.

\item The cycles of $s$, $r$, and $sr$ are as follows.  Parentheses denote
cycles on the coset labels, including fixed points:
\[
\begin{array}{c|l|l}
&50/\mathrm{II}&56/\mathrm{II}\\\hline
s&(01)(23)&(01)(24)(35)(69)(78)\\
r&(023)(1)&(023)(1)(467)(589)\\
sr&(012)(3)&(01265)(384)(79).
\end{array}
\]
Thus both curves have exactly one elliptic point of order $3$ and
none of order $2$.  Their cusp widths are $(3,1)$ and $(5,3,2)$.
The map $j$ to the $j$-line has respectively degree $4$ and $10$.
Its ramification contributions over $0,1728,\infty$ are
\[
 (2,2,2)\quad\hbox{and}\quad(6,5,7).
\]
Their sums are $2\cdot4-2$ and $2\cdot10-2$.  Riemann--Hurwitz
therefore gives genus zero in both cases.  In particular each base
is $\mathbb P^1$.

\item \par\medskip\noindent\textbf{2. The actual based peripheral loops.}\quad 
We describe a finite graph on each modular curve; this also specifies
the paths used for the matrices.  Put a black vertex at each
$r$-orbit and a white vertex at each $s$-orbit.  For every coset $q$,
join its black and white vertices by an edge labelled $q$, and
subdivide that edge at a midpoint also denoted by $q$.  The clockwise
cyclic order at a black vertex is the order of its $r$-cycle.  White
vertices have valence two.  This is the inverse image of an arc
joining the elliptic values on the modular sphere.  Filling its
complementary cusp discs recovers the compact curve.  Remove a
small disc about the fixed black vertex labelled $(1)$ as well.
The resulting surface is the base with its singular-fibre points
removed.  This description follows directly by lifting the modular
triangles: a three-cycle fills a disc above the order-three value,
an $s$-pair fills a disc above the order-two value, and an $sr$-cycle
is the boundary of a cusp disc.

\item The positive cusp boundary reads $sr=T$.  At the fixed black vertex,
the positive elliptic meridian reads $r^{-1}$.  These signs also
follow from the analytic coordinates: $\tau\mapsto\tau+1$ gives
a positive loop of $\exp(2\pi i\tau)$, while $R$ has derivative
$\exp(-2\pi i/3)$ at its fixed point in the upper half-plane.
We use the oriented elliptic marking in which this positive nodal
meridian acts by $T$; it fixes the choice between the two possible
outer identifications of the modular group with elliptic homology.

\item The base point is the midpoint $0$.  Stems to the midpoints $1,3,7$
are, respectively,
\[
 s,\qquad r^{-1},\qquad rsr^{-1};
\]
the last is needed only in the index-ten cover.  They are paths in
the graph just specified.  A cusp based at midpoint $q$, of width
$w$ and with stem $g$, therefore has based peripheral word
$gT^wg^{-1}$.  The fixed black vertex has word $sr^{-1}s^{-1}$.
Set
\begin{equation}\label{mapsD:repair:canonical-matrices}
 H=SR^{-1}S^{-1}=\begin{pmatrix}0&1\\-1&1\end{pmatrix},\quad
 V=R^{-1}TR=\begin{pmatrix}2&1\\-1&0\end{pmatrix},\quad
 B=V^3=\begin{pmatrix}4&3\\-3&-2\end{pmatrix},
\end{equation}
and
\[
 C=(RSR^{-1})T^2(RSR^{-1})^{-1}
   =\begin{pmatrix}5&2\\-8&-3\end{pmatrix}.
\]
The cusp and elliptic words, in positive boundary order, are
\[
 (T^3,V,H),\qquad (T^5,B,C,H).
\]

\item For completeness, we verify the boundary relation in the fundamental
group of the punctured surface, before imposing the order-three
orbifold relation.  Denote the black and white vertices by $b_i,w_i$
in the order in which their cycles occur in the displayed table.
Orient every half-edge from its midpoint $q$ towards its black or
white endpoint.  Take as a spanning tree every half-edge except
\[
\begin{array}{c|l}
50/\mathrm{II}&3\longrightarrow w_1,\\
56/\mathrm{II}&8\longrightarrow w_4,\quad9\longrightarrow w_3.
\end{array}
\]
The indicated subgraphs are connected and have one fewer edge than
vertices, hence are trees.  In the index-four case let $a$ denote
the remaining oriented half-edge, after the tree is contracted.
In the index-ten case call the two remaining half-edges $a,b$, in
the order displayed.  In both cases let $e$ denote the clockwise
loop at the punctured fixed black vertex.  These are free generators
of the actual punctured-surface group.  Reading the paths above and
discarding tree edges gives
\begin{equation}\label{mapsD:repair:free-peripherals}
\begin{array}{c|ccc}
50/\mathrm{II}&T^3&e a^{-1}&\\[-3pt]
&V&a&\\[-3pt]
&H&e^{-1}&
\end{array}
\qquad
\begin{array}{c|cc}
56/\mathrm{II}&T^5&e b^{-1}\\
&B&a\\
&C&a^{-1}b\\
&H&e^{-1}.
\end{array}
\end{equation}
For example the index-ten width-five cycle reads successively the
midpoints $0,1,2,6,5$; it crosses the punctured vertex once and the
half-edge $9\to w_3$ with negative orientation, giving $eb^{-1}$.
The width-three cycle $3,8,4$ gives $a$, and the width-two cycle
$7,9$ gives $a^{-1}b$.  In the index-four case the cycles $0,1,2$
and $3$ give $ea^{-1}$ and $a$.  A positive loop at the elliptic
puncture is $e^{-1}$ in both cases.  Thus the displayed boundary
products freely reduce to $1$ and their factors generate the whole
free group: in the second case $a=B$, $b=BC$, and $e=H^{-1}$.
The stems can be separated in a thin regular neighbourhood of the
tree, so they are genuine based peripheral loops, with the listed
cyclic order.  In particular this check does not replace the
punctured-surface group by its quotient in $G$.
\par\medskip\noindent\textbf{3. Weierstrass surfaces with these covers.}\quad 
For the index-four cover an explicit model is
\begin{equation}\label{mapsD:repair:50-weierstrass}
\begin{split}
 c_4&=t^2(t+27)(t+3),\\
 c_6&=t^3(t+27)(t^2+18t-27),\\
 \Delta&=t^7(t+27)^2,\\
 y^2&=x^3-\frac{c_4}{48}x-\frac{c_6}{864}.
\end{split}
\end{equation}
Multiplication gives $c_4^3-c_6^2=1728\Delta$.  The coefficients
are sections of $\mathcal O(4)$ and $\mathcal O(6)$ on
$\mathbb P^1$.  Their vanishing orders $(c_4,c_6,\Delta)$ at
$0,-27,\infty$ are, respectively,
\[
 (2,3,7),\qquad(1,1,2),\qquad(0,0,3).
\]
The minimal resolved fibres are therefore $I_1^*,\mathrm{II},I_3$.
No other fibre is singular.  Its $j$-map is
\[
 j(t)=\frac{(t+27)(t+3)^3}{t},\qquad
 j(t)-1728=\frac{(t^2+18t-27)^2}{t}.
\]
It has precisely the index-four passport in Step~1.  That passport
specifies its marked cover up to a relabelling of sheets: after
fixing $r=(023)(1)$, the three possible fixed-point-free involutions
on four letters are conjugate by the centralizer $\langle r\rangle$.
Thus the explicit model can be marked by the first coset table.

For the index-ten cover, the modular curve and its normalized
$j$-map specify the construction without the solution of any
additional polynomial equation.  Write $a$ for the single simple
point over $0$, $b_1,b_2,b_3$ for the triple points over $0$,
$d_1,\ldots,d_5$ for the double points over $1728$, and
$p_5,p_3,p_2$ for the three poles, whose subscripts are their orders.
All these points are distinct.  Choose a nonzero section
$\Delta$ of $\mathcal O(12)$ with divisor
\[
 5p_5+3p_3+2p_2+2a.
\]
The sections $j\Delta$ and $(j-1728)\Delta$ then have divisors
\[
 3(a+b_1+b_2+b_3),\qquad
 2(a+d_1+\cdots+d_5),
\]
respectively.  On $\mathbb P^1$ a section with either displayed
divisor has a cube root in $H^0(\mathcal O(4))$, respectively a
square root in $H^0(\mathcal O(6))$, after multiplying a section
with the divided divisor by a constant.  The required constant
roots exist over $\mathbb C$.  Choose these roots and call them
$c_4,c_6$.  They satisfy identically
\[
 c_4^3=j\Delta,\qquad c_6^2=(j-1728)\Delta,
 \qquad c_4^3-c_6^2=1728\Delta.
\]
The same Weierstrass equation as in
\eqref{mapsD:repair:50-weierstrass} therefore defines an elliptic
surface with a zero section.  At $a$ the orders are $(1,1,2)$,
giving type~II.  At $p_m$ they are $(0,0,m)$, giving $I_m$.
Every other point has nonvanishing discriminant.  The model is
minimal at each point since the orders of its two coefficients
are never simultaneously at least $4$ and $6$.  Its minimal
resolution is projective and has fundamental line bundle
$\mathcal O(1)$.  Both constructions are rational elliptic
surfaces by the minimal degree-$(4,6)$ criterion
\cite[\S8.2--8.4, pp.~38--39]{manSS10}.

\par\medskip\noindent\textbf{4. Lifting the projective marking.}\quad 
On a simply connected open set avoiding the singular values and
the elliptic values of $j$, uniformization identifies the fibre
with $\mathbb C/(\mathbb Z+\tau\mathbb Z)$, where $\tau$ is a
local inverse of the modular $j$-function.  Continuation through
the marked cover gives precisely its projective elliptic
monodromy.  The sign of an integral lift at each singular fibre
is fixed by its Kodaira type: $I_m$ has the positive unipotent
lift, $I_m^*$ its negative, and the type-II lift has trace~$1$.
With the peripheral paths of Step~2 the two ordered integral
tuples are consequently
\[
 (H,T^3,-V),\qquad (C,H,T^5,B).
\]
Their products are $I$, as direct multiplication verifies.  The
first tuple uses the negative cusp lift only at the width-one
cusp.  There is no residual quadratic character to choose: the
meridians generate the punctured-sphere group and their signs
have all been specified.  The cyclic rotations from Step~2
are geometric changes of the starting boundary arc.

Conjugating the first tuple simultaneously by $V$, and the second
by $B$, gives
\[
\begin{aligned}
 &(VHV^{-1},VT^3V^{-1},-V)\\[-1pt]
 &\quad=\left(
 \begin{pmatrix}3&7\\-1&-2\end{pmatrix},
 \begin{pmatrix}7&12\\-3&-5\end{pmatrix},
 \begin{pmatrix}-2&-1\\1&0\end{pmatrix}\right),\\
 &(BCB^{-1},BHB^{-1},BT^5B^{-1},B)\\[-1pt]
 &\quad=\left(
 \begin{pmatrix}5&8\\-2&-3\end{pmatrix},
 \begin{pmatrix}27&37\\-19&-26\end{pmatrix},
 \begin{pmatrix}61&80\\-45&-59\end{pmatrix},
 \begin{pmatrix}4&3\\-3&-2\end{pmatrix}\right).
\end{aligned}
\]
These are the stated matrices, in the stated order, in integral
oriented homology bases.

\par\medskip\noindent\textbf{5. The integral cocycle group and actual sections.}\quad 
We spell out the relation between
\eqref{mapsD:repair:cocycle-relation} and the geometric section
group, so that a rational height computation alone is not used
as an existence argument.  For any of the rational elliptic
surfaces just constructed, let $L$ be the orthogonal complement
of the unimodular sublattice generated by the fibre $F$ and zero section $O$ in integral $H_2$, with matrix $\left(\begin{smallmatrix}0&1\\1&-1\end{smallmatrix}\right)$ in the ordered basis $(F,O)$, and let $V_{\mathrm{fib}}$ be the subgroup generated by
the fibre components not meeting the zero section.  Fukae's
rational junction group $W_2$ consists of integral charges in
the rational local images, with total charge zero.  Its quotient
by integral null junctions $U_1$ is $L/V_{\mathrm{fib}}$;
its quotient by the saturation $U_2$ is the free Mordell--Weil
lattice, with the opposite pairing.  The quotient $U_2/U_1$
is the torsion subgroup
\cite[Lemma~4.1, Theorem~4.2 and Corollaries~4.2--4.3,
pp.~21--22]{manFukae02}.

Here is the exact convention conversion.  Put
\[
 y_i=P_iq_i,\qquad K_i=P_iA_i^{-1}P_i^{-1},
\]
and list the $K_i$ and $y_i$ in reverse order.  Then
$K_1\cdots K_n=I$, so the reversed list satisfies Fukae's
reverse-product convention.  His circuit matrices are inverses
of the ordinary monodromy matrices, as stated in his footnote
on p.~3.  The ordinary peripheral basis underlying the reversed
list is obtained by successive inverse Hurwitz moves:
its elements are $P_iA_iP_i^{-1}$, in reverse order.  This is a
geometric change of disjoint peripheral stems.  Moreover
\[
 y_i\in\operatorname{sat}\operatorname{im}(I-K_i),\qquad
 \sum_i y_i=0.
\]
Thus $(q_i)\mapsto(y_i)$ is an integral isomorphism from the
cocycle lattice to $W_2$.  Its inverse is $q_i=P_i^{-1}y_i$.

For an integral coboundary $q_i=(A_i-I)v$ we obtain
\[
 y_i=(P_{i+1}-P_i)v
     =(I-K_i)P_{i+1}v.
\]
In the reversed list, the product of the preceding Fukae
circuits is $P_{i+1}$.  The displayed formula is therefore
exactly his integral null-junction formula, and allowing
rational $v$ with integral charges gives its saturation
\cite[Lemmas~3.1--3.2, (11), pp.~11--13]{manFukae02}.
Consequently the quotient in
\eqref{mapsD:repair:cocycle-relation} is the full integral
section group, including its possible torsion.

The identification is the one for logarithm cocycles, not just
an abstract lattice isometry.  To see this, cut the punctured
base along its peripheral stems and trivialize the period
lattice on the cut disc.  A section has a continuous elliptic
logarithm there.  Along the two sides of the $i$th cut the
logarithms differ by $q_i$.  Transport to the common reference
fibre changes this charge to $P_iq_i=y_i$.  The boundary of
the swept relative chain is the sum of these charges, so the
closing condition is exactly $\sum_i y_i=0$.  Local rational
prongs fill precisely the rational image of $I-K_i$; retaining
integral charges gives $W_2$.  A change of the common logarithm
is the null sweep just calculated.  Conversely the junction
construction supplies the integral class in $L/V_{\mathrm{fib}}$
with these jumps.  On a rational surface the exponential
sequence identifies integral $H^2$ with the Picard group,
since $H^1(\mathcal O)=H^2(\mathcal O)=0$.  Restricting the
corresponding divisor class to the generic elliptic fibre
gives a degree-zero line bundle, hence its unique point under
the zero-section identification of the elliptic curve with
its degree-zero Picard variety.  That point extends to a
section of the minimal proper elliptic surface.  Vertical
components and the zero-section/fibre adjustment are exactly
the classes already quotiented out.  This reverses the
cut-and-jump construction and proves that every quotient
cocycle is realized by an actual section.

For profile $50/\mathrm{II}$ the first local difference is
\[
 A_1-I=\begin{pmatrix}2&7\\-1&-3\end{pmatrix},\qquad
 \det(A_1-I)=1.
\]
Every integral cocycle has a unique representative with
$q_1=0$, by subtracting one integral coboundary.  For the
displayed cocycle the required vector is $v=(1,-1)$, giving
\[
 q^0=\bigl((0,0),(-2,1),(2,-1)\bigr).
\]
The saturated image at $I_3$ is
$\mathbb Z(-2,1)$.  Once its coefficient $k\in\mathbb Z$
is chosen, the closing relation uniquely determines the
last charge, since $P_3$ is unimodular and the local image
at $I_1^*$ has full rank.  It is $k(2,-1)$.  Hence the
whole quotient is $\mathbb Zq^0$, with no torsion, and the
archived cocycle is a primitive section generator.

For profile $56/\mathrm{II}$ the type-II fibre is the second,
and again $\det(A_2-I)=1$.  Set $q_2=0$ uniquely by an
integral coboundary.  The other three saturated local images
have primitive generators
\[
 (-2,1),\qquad(4,-3),\qquad(1,-1).
\]
After prefix transport their columns are
\[
 (-2,1),\qquad(1,0),\qquad(1,-1).
\]
The first two have determinant $-1$.  The integral kernel
of the resulting $2$-by-$3$ matrix is exactly
$\mathbb Z(1,1,1)$, as the second row gives equality of the
first and third coefficients and the first row then gives
equality of the second.  Thus here too the quotient is
infinite cyclic, without torsion, and its generator is the
archived cocycle.

\par\medskip\noindent\textbf{6. Heights with the precise normalization.}\quad 
Let $z_i\in\mathbb Q^2$ satisfy $(A_i-I)z_i=q_i$ and write
$u\wedge v=u_1v_2-u_2v_1$.  In Fukae's convention a
preimage of $y_i$ under $I-K_i$ is
$P_i(z_i+q_i)$.  The diagonal contribution to the junction
square is therefore $z_i\wedge q_i$.  The reversal of the
ordering changes the cross term to
$-\sum_{i<j}(P_iq_i)\wedge(P_jq_j)$.  This gives
\begin{equation}\label{mapsD:repair:height-formula}
 \langle q,q\rangle_{\mathrm{MW}}
 =-\left\{\sum_i z_i\wedge q_i
       -\sum_{i<j}(P_iq_i)\wedge(P_jq_j)\right\}.
\end{equation}
For a multiplicative fibre this also reproduces Fukae's
$-t_i'^2/n_i$ term in (9): if
$A_i=I+n_i\,v(-v_2,v_1)$, substitution in $(A_i-I)z_i=q_i$
gives the same wedge expression.  Equivalently one may
check it on $A_i=T^{n_i}$, where $q_i=(t_i',0)$ and
$z_i=(0,t_i'/n_i)$, obtaining $-t_i'^2/n_i$.
The formula is independent of the choice of $z_i$, since
the kernel and image of a rank-one unipotent difference
are the same isotropic line.

For $50/\mathrm{II}$ take
\[
 (z_1,z_2,z_3)=
 \bigl((1,-1),(-4/3,0),(1/4,-3/4)\bigr).
\]
The prefix charges and the two sums in braces are
\phantomsection\label{qa:III:layout:2}
\[
\begin{gathered}
 (P_iq_i)=\bigl((-5,2),(4,0),(1,-2)\bigr),\\
 \sum_i z_i\wedge q_i=-3-\frac{16}{3}+\frac14
 =-\frac{97}{12},\qquad
 \sum_{i<j}(P_iq_i)\wedge(P_jq_j)=-8.
\end{gathered}
\]
Thus the junction square is $-1/12$ and the height is $1/12$.
For $56/\mathrm{II}$ take
\[
 (z_1,z_2,z_3,z_4)=
 \bigl((-1/2,0),(0,0),(-1/5,1/5),(1/3,0)\bigr).
\]
Here
\phantomsection\label{qa:III:layout:3}
\[
\begin{gathered}
 (P_iq_i)=\bigl((-2,1),(0,0),(1,0),(1,-1)\bigr),\\
 \sum_i z_i\wedge q_i=-\frac12-\frac15-\frac13
 =-\frac{31}{30},\qquad
 \sum_{i<j}(P_iq_i)\wedge(P_jq_j)=-1.
\end{gathered}
\]
The junction square is $-1/30$ and the height is $1/30$.
All entries in the theorem have now been realized with
the required based and integral markings.
\end{enumerate}
\end{proof}

\phantomsection\label{qa:III:layout:4}
\par\medskip\noindent\textbf{Exact finite verification.}\quad
The current verifier \path{code/III/receiving_final_certificates.py}
reconstructs both coset actions and the free peripheral words directly
from the punctured half-edge graphs. It compares both integral matrix
tuples and both cocycles with the archived marked profiles, checks the
Weierstrass identity, the primitive section quotients, the explicit
Fukae convention morphisms, and both rational heights. Its complete
finite output is \path{analysis/III_final_exact_certificates.json}.
The geometric existence and the actual logarithm-cocycle morphisms
are proved above; the program records the finite exact calculations.
